\documentclass[12pt,a4paper]{article}

% ==========================================
% Page layout
% ==========================================
\usepackage{geometry}
\geometry{
    top=2.0cm,
    bottom=2.0cm,
    left=2.1cm,
    right=2.1cm
}

% ==========================================
% XeLaTeX fonts
% ==========================================
\usepackage{fontspec}
\usepackage{unicode-math}

\setmainfont{TeX Gyre Termes}
\setmathfont{TeX Gyre Termes Math}

% ==========================================
% Packages
% ==========================================
\usepackage{amsmath}
\usepackage{graphicx}
\usepackage{float}
\usepackage{booktabs}
\usepackage{array}
\usepackage[hidelinks]{hyperref}
\usepackage{appendix}
\usepackage{setspace}

% Required line spacing
\onehalfspacing

% ==========================================
% Compact spacing
% ==========================================
\setlength{\parindent}{0pt}
\setlength{\parskip}{1pt}

\setlength{\textfloatsep}{4pt}
\setlength{\floatsep}{4pt}
\setlength{\intextsep}{4pt}

\setlength{\abovecaptionskip}{2pt}
\setlength{\belowcaptionskip}{0pt}

\setlength{\abovedisplayskip}{3pt}
\setlength{\belowdisplayskip}{3pt}
\setlength{\abovedisplayshortskip}{2pt}
\setlength{\belowdisplayshortskip}{2pt}

\begin{document}

% ==========================================
% Title
% ==========================================
\begin{center}
{\Large\bfseries
Statistical Analysis of the Average Velocity of a Cart on an Inclined Plane
}
\end{center}

% ==========================================
% Abstract
% ==========================================
\section*{Abstract}

The purpose of this experiment was to measure the average velocity of a cart moving down an inclined plane, evaluate the statistical uncertainty of repeated measurements, and compare the experimental result with an ideal kinematic model. The blocking time of a picket fence was measured over 150 trials using a photogate.

The experimental result was $\bar{v}{\mathrm{exp}}=(0.476\pm0.001),\mathrm{m/s}$, while the theoretical prediction was $v{\mathrm{theory}}=(0.548\pm0.008),\mathrm{m/s}$. The two values did not statistically agree. Friction, air resistance, and uncertainty in the effective release position were considered possible sources of the discrepancy.

% ==========================================
% Theoretical Background
% ==========================================
\section{Theoretical Background}

For a release distance \(s\) and photogate blocking length \(l\), the theoretical blocking time is

\begin{equation}
\Delta t=
\sqrt{\frac{2}{g\sin\theta}}
\left(\sqrt{s+l}-\sqrt{s}\right),
\label{eq:time}
\end{equation}

where \(g\) is the gravitational acceleration and \(\theta\) is the inclination angle. The theoretical average velocity while passing through the photogate is

\begin{equation}
v_{\mathrm{theory}}
=
\frac{l}{\Delta t}
=
\sqrt{\frac{g\sin\theta}{2}}
\left(\sqrt{s+l}+\sqrt{s}\right).
\label{eq:velocity}
\end{equation}

The model assumes that the cart is released from rest and that friction and air resistance are negligible. Detailed calculations are given in Appendix~\ref{app:theory}.

% ==========================================
% Materials and Methods
% ==========================================
\section{Materials and Methods}

A cart carrying a picket fence was released from rest on an inclined track toward a photogate. The photogate measured the blocking time \(\Delta t\) while the opaque section of the picket fence passed through it. The same release position and initial condition were used for 150 trials.

\begin{figure}[H]
    \centering
    \includegraphics[
        width=0.36\textwidth,
        trim=10 10 10 10,
        clip
    ]{exp1_syst.png}
    \caption{Experimental system showing the cart, inclined track and photogate.}
    \label{fig:system}
\end{figure}

The measured and given parameters were
\(g=(9.79127\pm0.000005)\,\mathrm{m/s^2}\),
\(l=(0.050\pm0.005)\,\mathrm{m}\),
\(s=(0.415\pm0.005)\,\mathrm{m}\), and
\(\theta=(2.00\pm0.05)^\circ\).
The measured blocking times were analyzed in Python, and each corresponding velocity was calculated from \(v_i=l/\Delta t_i\).

% ==========================================
% Results and Discussion
% ==========================================
\section{Results and Discussion}

\subsection{Theoretical and Experimental Results}

Using Eq.~\ref{eq:velocity}, the predicted velocity was
\(v_{\mathrm{theory}}=(0.548\pm0.008)\,\mathrm{m/s}\).
The uncertainty calculation is shown in Appendix~\ref{app:theory}.

A total of \(N=150\) blocking-time measurements were collected. Figure~\ref{fig:blocking_time} shows their distribution. Twelve bins were selected as a compromise: too few bins would hide important features of the distribution, while too many bins would make the histogram overly irregular because only a small number of measurements would fall into each bin.

\begin{figure}[H]
    \centering
    \includegraphics[
        width=0.44\textwidth,
        trim=15 10 15 10,
        clip
    ]{fig1.png}
    \caption{Distribution of the blocking time \(\Delta t\) for \(N=150\), with a normal fit.}
    \label{fig:blocking_time}
\end{figure}

The blocking time is centered at approximately \(0.1051\,\mathrm{s}\). With \(l=0.050\,\mathrm{m}\), this corresponds to a velocity of the order of \(0.5\,\mathrm{m/s}\), which is physically reasonable.

Each blocking time was converted into an average velocity. No outliers were identified using the \(3s\) criterion, so all 150 measurements were retained. The outlier calculation is shown in Appendix~\ref{app:statistics}. Figure~\ref{fig:velocity_dist} shows the resulting velocity distribution.

\begin{figure}[H]
    \centering
    \includegraphics[
        width=0.44\textwidth,
        trim=15 10 15 10,
        clip
    ]{fig2.png}
    \caption{Distribution of the experimental average velocities for \(N=150\), with the corresponding normal distribution.}
    \label{fig:velocity_dist}
\end{figure}

The histogram is approximately bell-shaped and follows the normal curve reasonably well. The experimental statistical results were
\(\bar{v}_{\mathrm{exp}}=0.476\,\mathrm{m/s}\),
\(s=0.0144\,\mathrm{m/s}\), and
\(s_m=0.0012\,\mathrm{m/s}\).
Thus, the experimental outcome is
\(\bar{v}_{\mathrm{exp}}=(0.476\pm0.001)\,\mathrm{m/s}\).
The calculation is shown in Appendix~\ref{app:statistics}.

The theoretical value lies more than two standard errors away from the experimental mean. Therefore, the experiment does not statistically agree with the ideal theoretical model. The statistical comparison is shown in Appendix~\ref{app:comparison}. The experimental value is approximately \(13\%\) lower than the theoretical prediction.

A possible source of this difference is friction between the cart and the track, since friction opposes relative motion between contacting surfaces [1]. Air resistance may also contribute [1]. Uncertainty in the effective release position may also affect the measured distance \(s\).

% ==========================================
% Central Limit Theorem
% ==========================================
\subsection{Effect of Sample Size and the Central Limit Theorem}

Subsets of \(N=20\) and \(N=60\) measurements were compared with the complete \(N=150\) dataset. Their statistical results are summarized in Table~\ref{tab:sample_analysis}.

\begin{table}[H]
    \centering
    \caption{Statistical parameters for different sample sizes.}
    \label{tab:sample_analysis}
    \begin{tabular}{cccc}
        \toprule
        \(N\) & Mean velocity (m/s) & SD (m/s) & SE (m/s) \\
        \midrule
        20  & 0.484 & 0.014 & 0.003 \\
        60  & 0.473 & 0.015 & 0.002 \\
        150 & 0.476 & 0.014 & 0.001 \\
        \bottomrule
    \end{tabular}
\end{table}

The standard deviations remain similar, while the standard error decreases as the sample size increases.

The velocity distribution for \(N=20\) is shown in Fig.~\ref{fig:n20}.

\begin{figure}[H]
    \centering
    \includegraphics[
        width=0.43\textwidth,
        trim=15 10 15 10,
        clip
    ]{fig3.png}
    \caption{Velocity distribution for \(N=20\), with the corresponding normal distribution.}
    \label{fig:n20}
\end{figure}

For \(N=20\), the histogram is relatively irregular and does not clearly show a normal shape. The corresponding \(N=60\) distribution is shown in Fig.~\ref{fig:n60}.

\begin{figure}[H]
    \centering
    \includegraphics[
        width=0.43\textwidth,
        trim=15 10 15 10,
        clip
    ]{fig4.png}
    \caption{Velocity distribution for \(N=60\), with the corresponding normal distribution.}
    \label{fig:n60}
\end{figure}

For \(N=60\), the distribution is smoother and follows the normal curve more closely. The complete \(N=150\) dataset gives the clearest normal-like distribution.

According to the Central Limit Theorem, approximately normal behavior is expected when the measurement is affected by many small random contributions. The results are consistent with this expectation: larger samples give a more stable distribution and a smaller standard error.

% ==========================================
% Conclusion
% ==========================================
\section{Conclusion}

The measured average velocity was
\((0.476\pm0.001)\,\mathrm{m/s}\), about \(13\%\) lower than the theoretical prediction
\((0.548\pm0.008)\,\mathrm{m/s}\). The two values do not statistically agree.

The difference may be related to friction and air resistance, which oppose the motion of the cart [1]. Uncertainty in the effective release position may also contribute. The repeated measurements were approximately consistent with a normal distribution, and the standard error decreased with increasing sample size.

% ==========================================
% References
% ==========================================
\section*{References}

[1] Urone, P. P., \& Hinrichs, R. (2022).
\textit{College Physics 2e}. OpenStax.
Chapter 5: Further Applications of Newton's Laws, Sections 5.1--5.2.
Retrieved from
\url{https://openstax.org/books/college-physics-2e/pages/5-section-summary}

% ==========================================
% Appendix
% ==========================================
\appendix
\section{Detailed Calculations}

{\small
\singlespacing

% ------------------------------------------
% A.1
% ------------------------------------------
\subsection{Theoretical Velocity and Uncertainty}
\label{app:theory}

For negligible friction, \(a=g\sin\theta\), and the velocity when the cart reaches the photogate is \(v_1=\sqrt{2gs\sin\theta}\). During the blocking interval,

\[
l=v_1\Delta t+\frac{1}{2}g\sin\theta(\Delta t)^2,
\]

which gives Eq.~\ref{eq:time}. Using Eq.~\ref{eq:velocity},

\[
v_{\mathrm{theory}}
=
\sqrt{\frac{g\sin\theta}{2}}
\left(\sqrt{s+l}+\sqrt{s}\right)
=
0.5481\,\mathrm{m/s}.
\]

The uncertainty was calculated using

\[
\delta v=
\sqrt{
\left(\frac{\partial v}{\partial g}\delta g\right)^2+
\left(\frac{\partial v}{\partial\theta}\delta\theta\right)^2+
\left(\frac{\partial v}{\partial l}\delta l\right)^2+
\left(\frac{\partial v}{\partial s}\delta s\right)^2
}.
\]

For Eq.~\ref{eq:velocity},

\[
\frac{\partial v}{\partial g}=\frac{v}{2g}, \qquad
\frac{\partial v}{\partial\theta}
=\frac{v}{2}\frac{\cos\theta}{\sin\theta},
\]

\[
\frac{\partial v}{\partial l}
=
\sqrt{\frac{g\sin\theta}{2}}
\frac{1}{2\sqrt{s+l}},
\qquad
\frac{\partial v}{\partial s}
=
\sqrt{\frac{g\sin\theta}{2}}
\left(
\frac{1}{2\sqrt{s+l}}+
\frac{1}{2\sqrt{s}}
\right).
\]

Using
\(\delta g=5\times10^{-6}\,\mathrm{m/s^2}\),
\(\delta l=\delta s=0.005\,\mathrm{m}\), and
\(\delta\theta=0.05^\circ=8.73\times10^{-4}\,\mathrm{rad}\),
the propagated uncertainty is
\(\delta v=0.0077\,\mathrm{m/s}\).

Thus,
\(v_{\mathrm{theory}}=(0.548\pm0.008)\,\mathrm{m/s}\).

% ------------------------------------------
% A.2
% ------------------------------------------
\subsection{Outlier and Statistical Calculations}
\label{app:statistics}

For the blocking-time data,

\[
\overline{\Delta t}
=
\frac{1}{N}\sum_{i=1}^{N}\Delta t_i,
\qquad
s_{\Delta t}
=
\sqrt{
\frac{
\sum_{i=1}^{N}
(\Delta t_i-\overline{\Delta t})^2
}{
N-1
}
}.
\]

Measurements outside
\(\overline{\Delta t}\pm3s_{\Delta t}\)
were considered potential outliers. No measurements were outside this interval.

For each measurement \(v_i=l/\Delta t_i\). The experimental mean, standard deviation and standard error were

\[
\bar v=\frac{1}{N}\sum v_i,
\qquad
s=
\sqrt{
\frac{\sum(v_i-\bar v)^2}{N-1}
},
\qquad
s_m=\frac{s}{\sqrt N}.
\]

The results were
\(\bar v=0.476\,\mathrm{m/s}\),
\(s=0.0144\,\mathrm{m/s}\), and
\(s_m=0.0012\,\mathrm{m/s}\).

% ------------------------------------------
% A.3
% ------------------------------------------
\subsection{Comparison with Theory}
\label{app:comparison}

According to the statistical criterion used in the course,

\[
\frac{|\bar v-v_{\mathrm{theory}}|}{s_m}\le2
\]

indicates reasonable agreement. For this experiment,

\[
\frac{|0.4763-0.5481|}{0.00118}
\approx61>2.
\]

Therefore, the experimental and theoretical results do not statistically agree.

}

\end{document}
